\Article{9}
Als zweites Beispiel wollen wir nun die Abbildung der Fläche eines geraden Kegels in der Ebene betrachten. Für die Kegelfläche nehmen wir die Gleichung
\[
xx+yy-kkzz=0
\]
\noindent an, wobei wir außerdem
\[
\begin{aligned}
x&=kt\cos u\\
y&=kt\sin u\\
z&=t
\end{aligned}
\]
\noindent und wie zuvor \(X=T,Y=U,Z=0\) setzen.

Die Differentialgleichung
\[
\omega=(kk+1)\mathrm{d}t^2+kktt\,\mathrm{d}u^2=0
\]
\noindent liefert hier die beiden Integrale
\[
\log t\pm i\BareRoot{\frac{kk}{kk+1}}\mathord{.}u=\text{Const.}
\]
\noindent Damit erhalten wir die Lösung
\[
\begin{aligned}
X+iY&=f(\log t+i\BareRoot{\frac{kk}{kk+1}}\mathord{.}u)\\
X-iY&=f'(\log t-i\BareRoot{\frac{kk}{kk+1}}\mathord{.}u)
\end{aligned}
\]
\noindent Das heißt, wenn \(f\) eine beliebige Funktion bezeichnet, nehmen wir für \(X\) den Realteil von
\[
f(\log t+i\BareRoot{\frac{kk}{kk+1}}\mathord{.}u)
\]
\noindent und für \(Y\) den Imaginärteil unter Weglassung des Faktors \(i\).
\par\indent\SourceLeaf{208}%
Setzt man für \(f\) zum Beispiel eine Exponentialgröße an, nämlich
\[
f\upsilon=he^{\upsilon}
\]
\noindent wobei \(h\) konstant ist und \(e\) die Basis der hyperbolischen Logarithmen bezeichnet, so erhält man die einfachste Abbildung
\[
X=ht\cos\BareRoot{\frac{kk}{kk+1}}\mathord{.}u,\quad
Y=ht\sin\BareRoot{\frac{kk}{kk+1}}\mathord{.}u
\]

Die Anwendung der Formeln aus Artikel 7 ergibt hier
\[
n=(kk+1)tt,\quad N=1
\]
\noindent und, da \(\varphi\upsilon=\varphi'\upsilon=he^{\upsilon}\),
\[
\varphi(\log t+i\BareRoot{\frac{kk}{kk+1}}\mathord{.}u)\mathord{.}\varphi'(\log t-i\BareRoot{\frac{kk'}{kk+1}}\mathord{.}u)=hhtt
\]
\noindent folglich
\[
m=\frac{h}{\BareRoot{(kk+1)}}
\]
\noindent also einen konstanten Wert. Setzt man außerdem
\[
h=\BareRoot{(kk+1)}
\]
\noindent so wird die Abbildung eine vollkommene Abwicklung.

